%% CAPSTONE Implementation Report
%% Powerline Detection with YOLinO + GNN
%% 코드 구현 명세 및 실험 환경 보고서
%%
%% Compile: pdflatex → bibtex → pdflatex × 2
%%
\documentclass[11pt,a4paper]{article}

%% ---- Packages ----
\usepackage[T1]{fontenc}
\usepackage[utf8]{inputenc}
\usepackage[english]{babel}
\usepackage{geometry}
\geometry{left=25mm, right=25mm, top=25mm, bottom=25mm}

\usepackage{amsmath, amssymb}
\usepackage{graphicx}
\usepackage{booktabs}
\usepackage{array}
\usepackage{multirow}
\usepackage{xcolor}
\usepackage{hyperref}
\hypersetup{colorlinks=true, linkcolor=blue, citecolor=blue, urlcolor=blue}
\usepackage{listings}
\usepackage{caption}
\usepackage{subcaption}
\usepackage{float}
\usepackage{enumitem}
\usepackage{fancyhdr}

%% ---- Code Listing Style ----
\definecolor{codegray}{rgb}{0.5,0.5,0.5}
\definecolor{codepurple}{rgb}{0.58,0,0.82}
\definecolor{backcolour}{rgb}{0.97,0.97,0.97}
\definecolor{keywordcolor}{rgb}{0.0,0.3,0.7}

\lstdefinestyle{pythonstyle}{
    backgroundcolor=\color{backcolour},
    commentstyle=\color{codegray}\itshape,
    keywordstyle=\color{keywordcolor}\bfseries,
    numberstyle=\tiny\color{codegray},
    stringstyle=\color{codepurple},
    basicstyle=\ttfamily\footnotesize,
    breakatwhitespace=false,
    breaklines=true,
    captionpos=b,
    keepspaces=true,
    numbers=left,
    numbersep=5pt,
    showspaces=false,
    showstringspaces=false,
    showtabs=false,
    tabsize=2,
    language=Python,
    frame=single,
    framerule=0.4pt,
    rulecolor=\color{black!30},
}
\lstset{style=pythonstyle}

%% ---- Header/Footer ----
\pagestyle{fancy}
\fancyhf{}
\lhead{\small CAPSTONE: Powerline Detection with YOLinO + GNN}
\rhead{\small Implementation Report}
\cfoot{\thepage}
\renewcommand{\headrulewidth}{0.4pt}

%% ============================================================
\begin{document}

\begin{titlepage}
  \centering
  \vspace*{2cm}
  {\LARGE \bfseries CAPSTONE: Powerline Detection with YOLinO + GNN\\[0.5em]}
  {\large \textit{Implementation Report --- Code Specification \& Experimental Environment}}\\[2em]
  \rule{\linewidth}{0.5pt}\\[1em]
  {\large Dataset, Preprocessing, Model Architecture, Training, Evaluation Metrics,\\
          Inference, and Ablation Studies}\\[1em]
  \rule{\linewidth}{0.5pt}\\[3em]
  {\large GitHub: \url{https://github.com/V8heart/CAPSTONE}}\\[0.5em]
  {\large Hugging Face: \url{https://huggingface.co/V8heart}}\\[3em]
  \vfill
  {\normalsize \today}
\end{titlepage}

\tableofcontents
\newpage

%% ============================================================
\section{Dataset and Preprocessing}
\label{sec:dataset}

%% ----
\subsection{TTPLA Dataset}
\label{subsec:ttpla}

The \textbf{TTPLA} (Transmission Tower and Power Line Aerial) dataset is the primary
benchmark used in this work. Its key statistics are as follows:

\begin{table}[H]
\centering
\caption{TTPLA dataset summary}
\begin{tabular}{ll}
\toprule
Property & Value \\
\midrule
Total images & 1,234 \\
Native resolution & 3840 $\times$ 2160 (4K UHD) \\
Instance annotations & 8,987 powerline instances \\
Annotation type & Polyline (vertex chain per wire) \\
Working resolution & 512 $\times$ 512 (tiled from 4K) \\
\bottomrule
\end{tabular}
\end{table}

The original 4K frames are preprocessed into $512 \times 512$ square crops via a
\emph{dual-crop} strategy (L/R or T/B depending on aspect ratio).
The processed benchmark is published on Hugging Face as
\texttt{V8heart/yolino-ttpla-benchmark} and is accessed via the
\texttt{DATASET\_TTPLA} environment variable.

Expected directory layout (set via \texttt{DATASET\_TTPLA}):
\begin{lstlisting}[language=bash, numbers=none, frame=none, backgroundcolor=\color{white}]
YOLinO_benchmark/
  images/{train,val,test}/*.png
  labels/{train,val,test}/*.npy     # polyline vertices per image
\end{lstlisting}

\subsubsection{1024$\times$1024 Ablation Tiles}
\label{subsec:ttpla_1024}

For backbone / predictor-count ablations (exp05--exp19) and full-resolution ISQ
evaluation, TTPLA is additionally tiled to $1024 \times 1024$ using the dual-crop
pipeline in \texttt{scripts/build\_ttpla\_style\_dataset.py}
(depends on \texttt{scripts/build\_downsampled\_yolino\_dataset.py}).
The processed dataset is published on Hugging Face as
\texttt{V8heart/ttpla-yolino-1024}.

\begin{table}[H]
\centering
\caption{TTPLA $1024 \times 1024$ tile split}
\begin{tabular}{lrr}
\toprule
Split & Images & Use \\
\midrule
train & 1,810 & Ablation training (50 epochs) \\
val   &   218 & Validation / FPS benchmark \\
test  &   440 & ISQ full-test evaluation \\
\bottomrule
\end{tabular}
\end{table}

\begin{lstlisting}[language=bash, numbers=none]
# Build 1024^2 tiles from CVAT XML (example)
python scripts/build_ttpla_style_dataset.py cvat-xml \
  --cvat-xml /path/to/annotations.xml \
  --image-dir /path/to/images \
  --guide-dir data_guide \
  --output-root ttpla_yolino_dataset_1024x1024 \
  --target-width 1024

# Download pre-built tiles
hf download V8heart/ttpla-yolino-1024 \
  --repo-type dataset \
  --local-dir ./ttpla_yolino_dataset_1024x1024
\end{lstlisting}

\subsubsection{Why TTPLA was Chosen}
Several candidate datasets were evaluated:

\begin{table}[H]
\centering
\caption{Candidate dataset comparison}
\begin{tabular}{lcccl}
\toprule
Dataset & Images & Resolution & Annotation & Decision \\
\midrule
\textbf{TTPLA} & 1,234 & 3840$\times$2160 & Polyline & \textbf{Selected} \\
Powerline IR/VL & $\sim$100 & $\leq$1080p & Bounding box & Insufficient scale \\
PLD-UAV & $\sim$9,000 & $\sim$1080p & Bounding box & No polyline GT \\
\bottomrule
\end{tabular}
\end{table}

TTPLA was selected because it provides (1) the largest publicly available 4K aerial
powerline dataset with per-instance \emph{polyline} annotations, (2) scene diversity
across 110+ locations, and (3) compatibility with the YOLinO grid regression framework.

\subsection{KEPCO Dataset}
\label{subsec:kepco}

In addition to TTPLA, a proprietary \textbf{KEPCO} (Korea Electric Power Corporation)
dataset is used for cross-domain evaluation (experiments exp71--exp75).

\subsubsection{Dataset Construction}
Raw footage from two field sources (minidata 001 and 002-part1) is
annotated in CVAT XML format and converted to the TTPLA-compatible
\texttt{\{images,labels\}/\{train,val,test\}} layout by
\texttt{scripts/prepare\_kepco\_yolino.py}
(uses \texttt{scripts/build\_downsampled\_yolino\_dataset.py}):

\begin{lstlisting}
# prepare_kepco_yolino.py (key pipeline)
crops = convert_cvat_dual_square_crops(
    polys_meta, img_path,
    target=target_width,   # 1024 px per side
    min_points=min_points
)
for tag, resized, polylines, instance_ids, tw, th in crops:
    stub = f"{b}_{tag}"          # e.g. "img001_L"
    resized.save(out_img)
    payload = {"polylines": polylines, "instance_ids": instance_ids}
    np.save(out_npy, payload)
\end{lstlisting}

The dual-crop strategy generates two $1024 \times 1024$ crops (left/right or top/bottom)
per source image, doubling effective dataset size while keeping annotations in the same
coordinate system as TTPLA.
The dataset split is guided by text files in \texttt{data\_guide/\{train,val,test\}.txt};
the final split is approximately 80:20 train/val.
The processed KEPCO benchmark is published on Hugging Face as
\texttt{V8heart/Yolino-KEPCO} with the same
\texttt{\{images,labels\}/\{train,val\}} layout used for exp23 (geom) and exp71 (GNN)
cross-domain evaluation.

\subsection{Label Format and Loading}
\label{subsec:label_loading}

Each \texttt{.npy} label supports two formats, both parsed in
\texttt{src/yolino/dataset/ttpla.py}:

\begin{lstlisting}
# ttpla.py -- TTPLADataset.__get_labels__()
raw = np.load(self.label_files[idx], allow_pickle=True)
payload = raw.item() if (isinstance(raw, np.ndarray)
                         and raw.dtype == object
                         and raw.shape == ()) else raw

if isinstance(payload, dict):
    # New format: {"instances": [{"instance_id":..., "points":...}, ...]}
    # or {"polylines": [...], "instance_ids": [...]}
    ...
else:
    # Legacy: raw list of polylines (fallback IDs 1..N)
    for p in payload:
        polylines.append(p)
        instance_ids.append(len(instance_ids) + 1)
\end{lstlisting}

Points are stored as $(x, y)$ in the file but converted to $(y, x)$ (row, col)
convention internally for the YOLinO grid pipeline. Each wire is sorted so that
$x_{\text{start}} \le x_{\text{end}}$ (left-to-right canonical order):

\begin{lstlisting}
# Canonical left-to-right ordering
if poly_np[0, 0] > poly_np[-1, 0]:   # x-coord of first vs last point
    poly_np = poly_np[::-1]
for j, pt in enumerate(poly_np):
    x, y = pt
    gridable_lines[i, j, 0] = float(y)   # row
    gridable_lines[i, j, 1] = float(x)   # col
\end{lstlisting}

\subsection{Z-Score Normalisation}
\label{subsec:normalization}

Images are normalised using \textbf{ImageNet} mean/std (ConvNeXt-Tiny pretrained values):

\[
\mu = [0.485,\; 0.456,\; 0.406], \quad
\sigma = [0.229,\; 0.224,\; 0.225]
\]

The normalisation is implemented as \texttt{torchvision.transforms.Normalize} and is
inserted automatically when \texttt{backbone\_pretrained: True} in the YAML config:

\begin{lstlisting}
# augmentation.py -- DatasetTransformer.__init__()
self.norm_mean = [0.485, 0.456, 0.406]
self.norm_std  = [0.229, 0.224, 0.225]
...
if getattr(args, "backbone_pretrained", False) \
        and Augmentation.NORM not in self.methods:
    self.methods.append(Augmentation.NORM)
    Log.info("Auto-enabled augmentation 'normalize' for "
             "pretrained ConvNeXt backbone (ImageNet mean/std).")
\end{lstlisting}

\subsection{Data Augmentation}
\label{subsec:augmentation}

Four augmentation types are applied during training, activated via the YAML key
\texttt{augment: crop,rotation,jitter,erasing} in
\texttt{exp80\_ttpla\_512512\_scale16.yaml}.
All augmentations are implemented in \texttt{src/yolino/dataset/augmentation.py}.

\begin{table}[H]
\centering
\caption{Data augmentation summary}
\begin{tabular}{llp{7cm}}
\toprule
Augmentation & Class & Parameters \\
\midrule
Random Crop &
  \texttt{RandomCropWithLabels} &
  Crop ratio from $[\texttt{crop\_range},\, 1.0]$;
  labels intersected with crop box via Shapely geometry.
  \texttt{crop\_range} controls minimum retained fraction. \\[4pt]
Random Erasing &
  \texttt{torchvision.transforms.RandomErasing} &
  $p=0.2$, scale $(0.02, 0.05)$, ratio $(0.15, 2.5)$, value=random. \\[4pt]
Random Rotation &
  \texttt{RandomRotationWithLabels} &
  Angle range $\pm$\texttt{rotation\_range} (radians$\to$degrees);
  points rotated via polar-coordinate transform around image centre. \\[4pt]
Color Jitter &
  \texttt{torchvision.transforms.ColorJitter} &
  brightness $0.5$, contrast $0.5$, saturation $0.5$, hue $0.2$. \\
\bottomrule
\end{tabular}
\end{table}

The crop/rotation transforms apply geometric transforms to label polylines
(not just images), ensuring pixel-accurate correspondences:

\begin{lstlisting}
# RandomRotationWithLabels.forward()
angle = RandomRotation.get_params(self.degrees)
image = F.rotate(image, angle)
for i_idx, instance in enumerate(uv_lines):
    for p_idx, point in enumerate(instance):
        point[0] -= h / 2;  point[1] -= w / 2
        pol = t_cart2pol(point)
        pol[1] += math.radians(angle)     # rotate angle
        point = t_pol2cart(pol)
        point[0] += h / 2;  point[1] += w / 2
        uv_lines[i_idx, p_idx] = point
params["rrotate_angle"] = angle
\end{lstlisting}

%% ============================================================
\section{Model Architecture}
\label{sec:model}

\subsection{Overall Pipeline}
\label{subsec:pipeline}

The system is a two-stage pipeline:

\begin{enumerate}[leftmargin=*]
  \item \textbf{Stage 1 (Geometry Head)} --- Input image $\to$ ConvNeXt-Tiny backbone $\to$
        FPN + PANet Neck $\to$ Geometry Head $\to$ per-cell line segment predictions.
  \item \textbf{Stage 2 (GAT Head)} --- Top-$K$ high-confidence segments form a graph;
        Graph Attention Network (GAT) predicts edge connectivity; connected components
        yield instance-level polylines.
\end{enumerate}

\subsection{Backbone: ConvNeXt-Tiny}
\label{subsec:backbone}

The backbone is \textbf{ConvNeXt-Tiny} (torchvision), loaded with
\texttt{ConvNeXt\_Tiny\_Weights.DEFAULT} (ImageNet-1K pretrained).
Feature extractor node mapping:

\begin{table}[H]
\centering
\caption{ConvNeXt-Tiny stage layout and FPN feature tap points}
\begin{tabular}{llll}
\toprule
Torchvision node & Stage & Stride & Channels \\
\midrule
\texttt{features.3} (stage2 blocks) & C2 & $\times 8$  & 192 \\
\texttt{features.5} (stage3 blocks) & C3 & $\times 16$ & 384 \\
\texttt{features.7} (stage4 blocks) & C4 & $\times 32$ & 768 \\
\bottomrule
\end{tabular}
\end{table}

\begin{lstlisting}
# yolino_net.py -- ConvNeXt_FPN_Backbone.__init__()
if _HAS_NEW_TV_API:
    weights = ConvNeXt_Tiny_Weights.DEFAULT if pretrained else None
    base = convnext_tiny(weights=weights)
else:
    base = convnext_tiny(pretrained=pretrained)
self.body = create_feature_extractor(
    base, return_nodes=OrderedDict([
        ("features.3", "C2"),
        ("features.5", "C3"),
        ("features.7", "C4"),
    ])
)
\end{lstlisting}

\subsection{Neck: FPN + PANet (Bottom-Up Path)}
\label{subsec:fpn}

The neck combines a top-down FPN with a PANet-style bottom-up path.
All FPN levels use 256 output channels, group normalisation (32 groups),
and Kaiming-initialised lateral and smooth convolutions.

\subsubsection{Top-Down FPN}

\[
M_4 = \text{Conv}_{1\times1}(C_4),\quad
M_3 = \text{Conv}_{1\times1}(C_3) + \text{Upsample}(M_4),\quad
M_2 = \text{Conv}_{1\times1}(C_2) + \text{Upsample}(M_3)
\]
\[
P_3 = \text{GN}(\text{Conv}_{3\times3}(M_3)),\quad
P_2 = \text{GN}(\text{Conv}_{3\times3}(M_2))
\]

\subsubsection{Bottom-Up Fusion (PANet)}

\[
N_3 = \text{GN}\!\left(\text{Conv}_{3\times3}\!\left(\text{stride2}(P_2) + M_3\right)\right),\quad
P_4^{\text{BU}} = \text{GN}\!\left(\text{Conv}_{3\times3}\!\left(\text{stride2}(N_3) + M_4\right)\right)
\]

\begin{lstlisting}
# yolino_net.py -- ConvNeXt_FPN_Backbone.forward()
feats = self.body(x)
c2, c3, c4 = feats["C2"], feats["C3"], feats["C4"]
m4 = self.lateral_c4(c4)
m3 = self.lateral_c3(c3) + self._upsample_to(m4, c3.shape[-2:])
m2 = self.lateral_c2(c2) + self._upsample_to(m3, c2.shape[-2:])
p3 = self.smooth_norm_p3(self.smooth_p3(m3))
p2 = self.smooth_norm_p2(self.smooth_p2(m2))
# Bottom-up path (PANet)
if self.bottom_up is not None:
    p4 = self.bottom_up(p2, m3, m4)
return {"P2": p2, "P3": p3, "P4": p4}
\end{lstlisting}

For the final model (exp80/exp83), the geometry head operates on
\textbf{P3} (stride $\times16$, scale=16) at resolution $32 \times 32$
for a $512 \times 512$ input.

\subsection{Geometry Head}
\label{subsec:geom_head}

\subsubsection{Grid Decomposition}

The input image is divided into a regular grid.
Each grid cell predicts multiple line segments independently:

\begin{table}[H]
\centering
\caption{Grid head hyperparameters (exp80)}
\begin{tabular}{ll}
\toprule
Parameter & Value \\
\midrule
Input resolution & $512 \times 512$ \\
Feature level & P3 (stride $\times 16$) \\
Cell size (scale) & $16 \times 16$ pixels (scale=16) \\
Grid shape & $32 \times 32$ cells \\
Predictors per cell & 4 \\
Training variables & \texttt{geom}, \texttt{confidence} \\
Activations & \texttt{linear} (geom), \texttt{sigmoid} (conf) \\
Confidence threshold & 0.7 \\
\bottomrule
\end{tabular}
\end{table}

\subsubsection{Predictor Count Ablation}

The baseline (exp01/Darknet) uses \textbf{8 predictors} per cell.
Ablation experiment exp19 reduces this to \textbf{4 predictors}:

\begin{lstlisting}
# exp19_fpn_bottomup_p4_num_predictors4_1024.yaml
# exp19: exp10 baseline + predictor 8 -> 4 ablation.
# Hypothesis: 4 predictors are sufficient to cover the max wires per cell.
num_predictors: 4
\end{lstlisting}

The 4-predictor setting is retained in all final experiments (exp80, exp83)
as it achieves equivalent performance with lower computation.

\subsubsection{Confidence Thresholding}

After the sigmoid activation, cells with confidence $< 0.7$ are suppressed:

\begin{lstlisting}
# exp80_ttpla_512512_scale16.yaml
confidence: 0.5   # training matching threshold
# exp83 inference threshold (isq_core.py)
GEOM_CONF = 0.7   # post-processing confidence gate
\end{lstlisting}

\subsection{GAT Head (Stage 2)}
\label{subsec:gat_head}

\subsubsection{Segment-to-Node Conversion}

High-confidence geometry cell predictions are treated as \emph{nodes} in a graph.
The top-$K$ ($K = 384$) segments (ranked by confidence) are selected:

\begin{lstlisting}
# exp83_gnn_ttpla_512512_from_exp80.yaml
gnn_max_nodes: 384          # top-K segment cap
gnn_node_conf_thresh: 0.1   # pre-NMS node gate
\end{lstlisting}

Each node is initialised from its mid-direction geometry token
augmented with sinusoidal positional encoding and (optionally) visual FPN features.

\subsubsection{Edge Attention (GAT)}

The GAT head operates on the \textbf{directional2} adjacency mode:
for each node, at most $K_{\text{dir}} = 8$ directed neighbours
(forward/backward along segment direction) are considered as edge candidates.
Geometric constraints filter edges:

\begin{table}[H]
\centering
\caption{Geometric edge prior hyperparameters (exp83)}
\begin{tabular}{ll}
\toprule
Parameter & Value \\
\midrule
Max lateral distance & 5.0 px (symmetric) \\
Max along-line gap & 128.0 px \\
Max endpoint gap & 40.0 px \\
Min direction dot product & 0.96 \\
Min separation & 8.0 px \\
Directional neighbours $K$ & 8 \\
\bottomrule
\end{tabular}
\end{table}

The attention weight $e_{ij}$ between node $i$ and neighbour $j$ is computed as:

\[
e_{ij} = \frac{1}{\sqrt{d_h}} \mathbf{q}_i \cdot \mathbf{k}_j + b_{ij}^{\text{geom}}
\]

where $\mathbf{q}_i, \mathbf{k}_j \in \mathbb{R}^{d_h}$ are query/key projections from
$h = 4$ attention heads, and $b_{ij}^{\text{geom}}$ is an optional geometric bias.
Weights are softmax-normalised over the $K$ neighbour axis:

\begin{lstlisting}
# yolino_gnn_head.py -- _GATLayer.forward()
att = (q.unsqueeze(2) * k_neigh).sum(dim=-1) / math.sqrt(dh)  # [B,N,K,H]
if geom_bias is not None:
    att = att + geom_bias.unsqueeze(-1)
att = att.masked_fill(~neigh_valid.unsqueeze(-1), float("-inf"))
w = F.softmax(att_safe, dim=2)    # softmax over K neighbours
msg = (w.unsqueeze(-1) * v_neigh).sum(dim=2).reshape(b, n, hd)
x = x + self.resid_drop(self.proj_out(msg))
x = x + self.ffn(self.norm2(x))
\end{lstlisting}

\subsubsection{Edge Classification and Sigmoid Threshold}

An MLP edge head takes the concatenation $[\mathbf{h}_i, \mathbf{h}_j,
\mathbf{f}_{ij}^{\text{geom}}]$ and outputs a scalar logit per candidate edge.
After sigmoid, edges with probability $\geq 0.3$ (inference, \texttt{gnn\_overlay\_edge\_thresh})
are retained:

\begin{lstlisting}
# exp83_gnn_ttpla_512512_from_exp80.yaml
gnn_cc_edge_thresh: 0.2          # connected-component threshold (lowered from 0.3)
gnn_overlay_edge_thresh: 0.3     # visualisation/inference threshold
gnn_eval_edge_threshs: "0.2,0.35,0.5"
\end{lstlisting}

\subsubsection{90\% Masking and Top-384 Segment Matrix}

To limit quadratic memory/computation in the edge matrix, only the top-384
high-confidence segments are used as nodes. Edges between all $384 \times 384$
pairs are \emph{not} computed: the \texttt{directional2} mode restricts each node
to at most $K=8$ directional candidates, yielding a $384 \times 8$ sparse
edge matrix. Approximately 90\% of potential pairs are masked via geometric
hard-gating before attention:

\begin{lstlisting}
# segment_geom_soft_gate.py / yolino_gnn_head.py
# Hard gate: lateral and along constraints
gnn_use_hard_geom_gate: true
gnn_max_lateral_px: 5.0
gnn_max_lateral_sym: true
gnn_max_along_px: 128.0
gnn_max_end_gap_px: 40.0
# Result: ~90% of N x N pairs excluded before GAT attention
\end{lstlisting}

%% ============================================================
\section{Training}
\label{sec:training}

\subsection{Weight Initialisation}
\label{subsec:init}

All FPN lateral/smooth convolutions and the bottom-up path convolutions
use \textbf{Kaiming (MSRA) normal initialisation} (fan-out mode):

\begin{lstlisting}
# yolino_net.py -- ConvNeXt_FPN_Backbone.__init__()
for m in (self.lateral_c2, self.lateral_c3, self.lateral_c4,
          self.smooth_p2, self.smooth_p3, self.smooth_p4):
    nn.init.kaiming_normal_(m.weight,
                            mode="fan_out",
                            nonlinearity="relu")
    if m.bias is not None:
        nn.init.zeros_(m.bias)
\end{lstlisting}

The ConvNeXt-Tiny body itself retains ImageNet pretrained weights.

\subsection{Learning Rate Scheduler}
\label{subsec:scheduler}

A \textbf{warm-up cosine annealing} schedule is used in Stage 2 (exp83):

\[
\eta(t) =
\begin{cases}
\dfrac{t+1}{T_{\text{warm}}} \cdot \eta_0 & t < T_{\text{warm}} \\[6pt]
\eta_{\min} + \left(1 - \eta_{\min}/\eta_0\right)
  \cdot \dfrac{1 + \cos\!\left(\pi\,\dfrac{t - T_{\text{warm}}}{T - T_{\text{warm}}}\right)}{2}
& t \geq T_{\text{warm}}
\end{cases}
\]

where $\eta_0$ is the initial learning rate, $T$ is total training steps,
$T_{\text{warm}}$ is the number of warmup steps, and
$\eta_{\min} = 0.1 \cdot \eta_0$.

\begin{lstlisting}
# optimizer_factory.py -- get_optimizer()
def lr_lambda(step: int):
    if scheduler == Scheduler.WARMUP_COSINE and step < warmup_steps:
        return float(step + 1) / float(warmup_steps)
    cosine_total = max(1, total_steps - warmup_steps)
    cosine_step  = max(0, min(step - warmup_steps, cosine_total))
    cos_factor = 0.5 * (1.0 + math.cos(
                    math.pi * cosine_step / cosine_total))
    return min_lr_ratio + (1.0 - min_lr_ratio) * cos_factor
scheduler = torch.optim.lr_scheduler.LambdaLR(
    optimizer, lr_lambda=lr_lambda)
\end{lstlisting}

\begin{table}[H]
\centering
\caption{Scheduler hyperparameters (exp83)}
\begin{tabular}{ll}
\toprule
Parameter & Value \\
\midrule
Optimiser & Adam \\
GNN learning rate (\texttt{lr\_e2e}) & $1 \times 10^{-4}$ \\
Backbone/FPN/Geom LR & 0.0 (frozen) \\
Scheduler & \texttt{warmup\_cosine} (step-level) \\
Warmup epochs & 0 \\
Min LR ratio & 0.1 \\
Batch size & 16 \\
Total epochs & 50 \\
Early stopping patience & 25 epochs (earliest: 120) \\
Gradient clip norm & 1.0 \\
\bottomrule
\end{tabular}
\end{table}

\subsection{Fine-Tuning Strategy}
\label{subsec:finetuning}

\textbf{Stage 1 (exp80)}: All layers (ConvNeXt body + FPN + geometry head) are
trained end-to-end on TTPLA $512\times512$ from ImageNet pretrained weights.

\begin{lstlisting}
# exp80_ttpla_512512_scale16.yaml
backbone_pretrained: True       # load ImageNet weights
# No lr_backbone/lr_fpn overrides -> flat LR for all groups
learning_rate: (default, ~1e-3)
\end{lstlisting}

\textbf{Stage 2 (exp83)}: The Stage 1 checkpoint is loaded as \texttt{explicit\_model}.
Backbone, FPN, and geometry head are \emph{frozen} (learning rate = 0.0);
only the GNN head is trained:

\begin{lstlisting}
# exp83_gnn_ttpla_512512_from_exp80.yaml
explicit_model: "log/checkpoints/exp80_ttpla_512512_scale16/best_model.pth"
lr_backbone: 0.0
lr_fpn:      0.0
lr_geom:     0.0
lr_e2e:      0.0001    # GNN head only
backbone_freeze_epochs: 10
e2e_mode: gnn
\end{lstlisting}

\begin{lstlisting}
# optimizer_factory.py -- maybe_freeze_backbone()
def maybe_freeze_backbone(args, net, epoch):
    n_freeze = int(getattr(args, "backbone_freeze_epochs", 0) or 0)
    if n_freeze <= 0:
        return False
    should_freeze = epoch < n_freeze
    for name, p in net.named_parameters():
        if any(name.startswith(p) for p in _BACKBONE_BODY_PREFIXES):
            p.requires_grad = not should_freeze
\end{lstlisting}

\subsection{Loss Functions}
\label{subsec:loss}

\subsubsection{Stage 1 --- Geometry + Confidence Loss}

Stage 1 minimises a weighted combination of:
\begin{itemize}
  \item \textbf{Geometry loss} (\texttt{norm\_mean}): $L_2$ normalised MSE on
        segment endpoint / mid-dir parameters.
  \item \textbf{Confidence loss} (\texttt{mse\_mean}): MSE between predicted and
        binary GT confidence (1 for matched cells, 0 for empty cells).
\end{itemize}

\[
\mathcal{L}_{\text{Stage1}} = 2 \cdot \mathcal{L}_{\text{geom}} + 1 \cdot \mathcal{L}_{\text{conf}}
\]

Loss weight strategy is \texttt{fixed} with \texttt{weights: "2,1"}.
Hard-matching bipartite assignment is used (\texttt{loss\_hard\_matching: true},
\texttt{match\_by\_conf\_first: true}).

\subsubsection{Stage 2 --- GNN Edge Loss}

Stage 2 adds a \textbf{Binary Cross-Entropy} (BCE) edge classification loss
with positive-class reweighting:

\[
\mathcal{L}_{\text{edge}} = -\frac{1}{|\mathcal{E}|}
\sum_{(i,j)\in\mathcal{E}} \left[
  w^+ y_{ij} \log \sigma(\hat{e}_{ij}) + (1 - y_{ij}) \log(1 - \sigma(\hat{e}_{ij}))
\right]
\]

\begin{lstlisting}
# exp83_gnn_ttpla_512512_from_exp80.yaml
gnn_edge_loss_type: bce
gnn_pos_weight: 4.0          # positive edge weight (2.0 in exp82 -> 4.0 in exp83)
gnn_neg_per_pos: 6           # hard-negative ratio  (8 in exp82 -> 6 in exp83)
gnn_neg_min_kept: 2
gnn_loss_weight: 0.05        # combined GNN loss scale
\end{lstlisting}

An additional \textbf{random-walk topology loss} regularises global connectivity
over $k=8$ steps (increased from 6 in exp82 to force longer-range chain connectivity):

\[
\mathcal{L}_{\text{RW}} = \text{BCE}\!\left(T^k_{ij},\;
\mathbf{1}[\text{same instance}_{ij}]\right), \quad k=8 \text{ steps}
\]

\begin{lstlisting}
# gnn_topology_loss.py
# exp83: steps=8, rw_pos_weight=30.0, rw_topology_weight=0.35
tk = transition
for _ in range(steps - 1):      # T^k via matrix multiply
    tk = torch.bmm(tk, transition)
bce = F.binary_cross_entropy(pred, target, reduction="none")
loss = (bce * weight)[valid_pair].sum() / weight[valid_pair].sum()
\end{lstlisting}

\begin{table}[H]
\centering
\caption{GNN edge loss hyperparameters: exp82 vs.\ exp83}
\begin{tabular}{llll}
\toprule
Parameter & exp82 & \textbf{exp83} & Effect \\
\midrule
\texttt{gnn\_pos\_weight}          & 2.0 & \textbf{4.0}  & Stronger positive edge signal \\
\texttt{gnn\_neg\_per\_pos}        & 8   & \textbf{6}    & Compensates higher pos weight \\
\texttt{gnn\_degree\_row/col\_cap} & 2.0 & \textbf{4.0}  & Allows bridge edges \\
\texttt{gnn\_rw\_topology\_weight} & 0.1 & \textbf{0.35} & Stronger topology regularisation \\
\texttt{gnn\_rw\_steps}            & 6   & \textbf{8}    & Longer chain connectivity \\
\texttt{gnn\_rw\_pos\_weight}      & 20.0& \textbf{30.0} & Same-instance pairs amplified \\
\texttt{gnn\_rw\_start\_epoch}     & 3   & \textbf{2}    & RW loss starts 1 epoch earlier \\
\texttt{gnn\_cc\_edge\_thresh}     & 0.3 & \textbf{0.2}  & More segments connected at CC \\
\texttt{gnn\_overlay\_edge\_thresh}& 0.35& \textbf{0.3}  & Consistent with CC threshold \\
\bottomrule
\end{tabular}
\end{table}

%% ============================================================
\section{Baseline and Ablation Studies}
\label{sec:ablation}

\subsection{Ablation Design}

Three ablation steps progressively upgrade the system:

\begin{table}[H]
\centering
\caption{Ablation experiment configurations (1024$\times$1024 TTPLA, unless noted)}
\begin{tabular}{llllll}
\toprule
Exp & Config key & Backbone & Neck & Predictors & Head level \\
\midrule
exp09 & \texttt{exp09\_darknet\_1024} & Darknet-19 & None & 8 & P4 (stride 32) \\
exp05 & \texttt{exp05\_convnext\_baseline\_1024} & ConvNeXt-Tiny & None & 8 & P4 \\
exp10 & \texttt{exp10\_fpn\_bottomup\_p4\_1024} & ConvNeXt-Tiny & FPN+BU & 8 & P4 \\
exp19 & \texttt{exp19\_fpn\_bottomup\_p4\_num\_predictors4\_1024} & ConvNeXt-Tiny & FPN+BU & 4 & P4 \\
\textbf{exp80} & \texttt{exp80\_ttpla\_512512\_scale16} & ConvNeXt-Tiny & FPN+BU & 4 & P3 (stride 16) \\
\textbf{exp83} & \texttt{exp83\_gnn\_ttpla\_512512\_from\_exp80} & (frozen) & (frozen) & 4 & P3 + GNN \\
\bottomrule
\end{tabular}
\end{table}

\subsection{Ablation Results (validation F1)}

Each ablation row was trained for 50 epochs on TTPLA at $1024 \times 1024$.
The reported F1 is the \textbf{best validation F1} logged during training
(cell-level matching metric from the YOLinO training loop, not LSNetv2 pixel F1).
Values are taken from
\texttt{ttpla\_train\_exp/analysis/metric\_histograms/best\_f1\_epoch\_metrics\_summary.csv}.

\begin{table}[H]
\centering
\caption{Ablation study: progressive architecture upgrades (slide~18)}
\begin{tabular}{clrrrr}
\toprule
Step & Experiment & Best epoch & F1 & Precision & Recall \\
\midrule
1 & exp09 (Darknet-19 baseline)           & 28 & 0.677 & 0.585 & 0.808 \\
2 & exp05 (+ ConvNeXt-Tiny backbone)        & 36 & 0.748 & 0.637 & 0.911 \\
3 & exp10 (+ top-down FPN + bottom-up)      & 29 & 0.784 & 0.689 & 0.915 \\
4 & exp19 (+ reduce predictors $8 \to 4$) & 41 & 0.836 & 0.782 & 0.901 \\
\bottomrule
\end{tabular}
\end{table}

\noindent
Reproduce by training each YAML under \texttt{configs/experiments/} and
inspecting TensorBoard / the CSV above.
The production Stage~1 model for $512^2$ inference is \textbf{exp80}
(further fine-tuned from the exp19 family at P3 / scale~16).

\subsection{Toggling via Config Files}

Each ablation is a separate YAML; backbone, FPN, and predictor count are
toggled via top-level keys:

\begin{lstlisting}
# exp09 (baseline): Darknet, no FPN
backbone: darknet
use_fpn: false
num_predictors: 8
head_level: P4

# exp05: swap backbone only
backbone: convnext
use_fpn: false
num_predictors: 8

# exp10: add FPN + bottom-up
backbone: convnext
use_fpn: true
use_bottom_up: true
num_predictors: 8

# exp19: reduce predictors 8->4
backbone: convnext
use_fpn: true
use_bottom_up: true
num_predictors: 4   # key change

# exp80: move head to P3 (scale=16)
scale: 16
head_level: P3
num_predictors: 4

# exp83: Stage 2 GNN on top of exp80
#  - backbone/fpn/geom frozen (lr=0.0)
#  - GNN trained with strengthened topology loss
e2e_mode: gnn
explicit_model: "log/checkpoints/exp80_ttpla_512512_scale16/best_model.pth"
lr_backbone: 0.0;  lr_fpn: 0.0;  lr_geom: 0.0
lr_e2e: 0.0001
gnn_rw_topology_weight: 0.35   # 0.1 in exp82
gnn_cc_edge_thresh: 0.2        # 0.3 in exp82
\end{lstlisting}

At runtime, \texttt{YolinoNet} selects the backbone class
(\texttt{ConvNeXt\_FPN\_Backbone} vs.\ \texttt{Darknet\_FPN\_Backbone})
and configures the FPN and head level:

\begin{lstlisting}
# yolino_net.py -- YolinoNet.__init__()
pretrained = bool(getattr(args, "backbone_pretrained", True))
if args.backbone == "convnext":
    self.backbone = ConvNeXt_FPN_Backbone(
        out_channels=fpn_out,
        pretrained=pretrained,
        use_fpn=args.use_fpn,
        use_bottom_up=args.use_bottom_up,
    )
elif args.backbone == "darknet":
    self.backbone = Darknet_FPN_Backbone(
        pretrained=pretrained,
        use_fpn=args.use_fpn,
        ...
    )
self.head_level = str(getattr(args, "head_level", "P3"))
\end{lstlisting}

%% ============================================================
\section{Evaluation Metrics}
\label{sec:metrics}

\subsection{Pixel-Level Metrics (APR / ARR / F1 / F$_\beta$)}
\label{subsec:pixel_metrics}

\subsubsection{Rasterisation}

Both GT polylines and predicted polylines/segments are rasterised to binary masks
using \texttt{cv2.line} at a specified line width.
The LSNetv2 protocol uses widths of \textbf{2 px} and \textbf{4 px}:

\begin{lstlisting}
# eval_pixel_f1/eval_pixel_f1.py
LSNETV2_LINE_WIDTHS = (2, 4)   # px
LSNETV2_BETA_SQ = 0.3          # beta^2 for F_beta

def rasterize_polylines_mask(height, width, polylines, thickness):
    return build_binary_mask(height, width, polylines, thickness)

def rasterize_line_segments_mask(height, width, segments,
                                  thickness):
    mask = np.zeros((height, width), dtype=np.uint8)
    for seg in segments:
        (x1, y1), (x2, y2) = seg
        cv2.line(mask,
                 (int(round(x1)), int(round(y1))),
                 (int(round(x2)), int(round(y2))),
                 1,
                 max(1, int(thickness)))
    return mask
\end{lstlisting}

\subsubsection{Per-Image Precision / Recall}

For each image, pixel-level true positives (TP), false positives (FP),
and false negatives (FN) are computed from the binary mask intersection:

\begin{lstlisting}
# per-image computation
tp = int((pred_mask & gt_mask).sum())
fp = int((pred_mask & ~gt_mask).sum())
fn = int((~pred_mask & gt_mask).sum())
precision = tp / (tp + fp + 1e-8)
recall    = tp / (tp + fn + 1e-8)
f1        = 2*precision*recall / (precision + recall + 1e-8)
\end{lstlisting}

\subsubsection{Macro APR / ARR / F1 / F$_\beta$}

Following the LSNetv2 protocol, \emph{macro} averages are reported
(average over images, not global pixel count):

\[
\mathrm{APR} = \frac{1}{N}\sum_{i=1}^N P_i,\quad
\mathrm{ARR} = \frac{1}{N}\sum_{i=1}^N R_i
\]
\[
F_1 = \frac{2 \cdot \mathrm{APR} \cdot \mathrm{ARR}}
           {\mathrm{APR} + \mathrm{ARR}},\quad
F_\beta = \frac{(1+\beta^2)\,\mathrm{APR}\cdot\mathrm{ARR}}
               {\beta^2\,\mathrm{APR} + \mathrm{ARR}}, \quad \beta^2 = 0.3
\]

\begin{lstlisting}
# eval_pixel_f1.py
def f_beta_score(precision, recall, beta_sq=LSNETV2_BETA_SQ):
    return float((1.0 + beta_sq) * precision * recall
                 / (beta_sq * precision + recall + 1e-8))

def compute_lsnetv2_metrics(per_image):
    apr = float(np.mean([m["precision"] for m in per_image]))
    arr = float(np.mean([m["recall"]    for m in per_image]))
    f1  = 2*apr*arr / (apr + arr + 1e-8)
    return {"APR": apr, "ARR": arr,
            "F1": float(f1), "F_beta": f_beta_score(apr, arr)}
\end{lstlisting}

\begin{table}[H]
\centering
\caption{Pixel F1 results on TTPLA test set (220 images, 2 px line width, GT-B)}
\begin{tabular}{llrrrr}
\toprule
Model & Stage & APR & ARR & F1 & F$_\beta$ ($\beta^2{=}0.3$) \\
\midrule
exp80 & Geom only & 0.610 & 0.581 & 0.595 & 0.603 \\
exp83 (ep6) & + GNN & 0.617 & 0.579 & 0.597 & 0.608 \\
\midrule
\multicolumn{2}{l}{LSNetv2 (reported, 512$\times$512)} & 0.714 & 0.560 & 0.628 & 0.671 \\
\bottomrule
\end{tabular}
\end{table}

\subsection{ISQ: Instance Segmentation Quality}
\label{subsec:isq}

ISQ is a proposed segment-level metric that accounts for instance identity
(polyline ID), over-splitting, and under-merging.
All implementation is in \texttt{eval\_isq/isq\_core.py}.

\subsubsection{Constants}

\begin{lstlisting}
# isq_core.py
CELL_SIZE = 32
MATCH_MAX_DIST_PX  = 24.0    # distance threshold
MATCH_MAX_ANGLE_DEG = 15.0   # angle threshold
MEANINGFUL_COVERAGE_THRESH = 0.3   # 30% overlap for UM/OS
MIN_SEGMENTS_PER_PRED_POLYLINE = 2
\end{lstlisting}

\subsubsection{Step 1: GT Segment Generation}

Each GT polyline is intersected with a $32 \times 32$ pixel grid.
At most one segment per (cell, polyline) pair is generated:

\begin{lstlisting}
# isq_core.py -- polyline_to_cell_segments()
for row in range(n_rows):
    for col in range(n_cols):
        seg = polyline_segment_in_cell(
            polyline, row, col, cell_size=CELL_SIZE)
        if seg is None:
            continue
        (x1, y1), (x2, y2) = seg
        out.append(GtSegment(x1=x1, y1=y1, x2=x2, y2=y2,
                             polyline_id=int(polyline_id),
                             seg_idx=idx))
        idx += 1
\end{lstlisting}

\subsubsection{Step 2: Predicted Segment to GT Matching}

For each predicted segment, the closest GT segment is found subject to
distance $< 24\,\text{px}$ \textbf{AND} angle difference $< 15°$:

\[
d(s_p, s_g) = \max\bigl(
  d_\perp(A_p, s_g),\; d_\perp(B_p, s_g),\;
  d_\perp(A_g, s_p),\; d_\perp(B_g, s_p)\bigr)
\]

\begin{lstlisting}
# isq_core.py -- match_pred_segment_to_gt()
for gt in candidates:
    gt_seg = gt.as_endpoints()
    dist  = segment_distance(pred_seg, gt_seg)
    if dist >= MATCH_MAX_DIST_PX:        continue
    angle = angle_diff_deg(segment_angle(pred_seg),
                           segment_angle(gt_seg))
    if angle >= MATCH_MAX_ANGLE_DEG:     continue
    if dist < best_dist:
        best_dist = dist;  best_idx = gt.seg_idx
\end{lstlisting}

A spatial bucket index (\texttt{build\_gt\_spatial\_index}) accelerates candidate
retrieval from $O(N_{\text{gt}})$ to $O(1)$ per query.

\subsubsection{Step 3: Dominant-ID and TP/FP/FN Computation}

For each predicted polyline $p_i$ (with $\geq 2$ segments), the
\emph{dominant GT polyline} is the most-frequently matched GT instance:

\begin{lstlisting}
# isq_core.py -- _dominant_gt_id_for_pred()
ids = [m.matched_polyline_id
       for m in matches
       if m.pred_polyline_idx == pi
       and m.matched_polyline_id is not None]
dominant_id = Counter(ids).most_common(1)[0][0]
\end{lstlisting}

TP/FP/FN are computed ensuring \emph{each GT segment is counted as TP at most once}:

\begin{lstlisting}
# isq_core.py -- compute_isq_image()
claimed_gt: set[int] = set()
tp = fp = 0
for m in matches:
    dom = dominant_ids.get(m.pred_polyline_idx)
    if dom is None or m.matched_polyline_id != dom \
            or m.matched_gt_seg_idx is None:
        fp += 1;  continue
    if m.matched_gt_seg_idx in claimed_gt:
        fp += 1;  continue      # GT already claimed
    claimed_gt.add(m.matched_gt_seg_idx)
    tp += 1
fn = sum(1 for g in gt_segments
         if g.seg_idx not in claimed_gt)
\end{lstlisting}

\subsubsection{Over-Split (OS) Rate}

An over-split occurs when $\geq 2$ predicted polylines each have
$\geq 30\%$ coverage of the same GT polyline:

\begin{lstlisting}
# isq_core.py -- _count_over_split_instances()
for gid in gt_ids:
    n_gt = len(gt_seg_by_id[gid])
    meaningful_preds = 0
    for pi in range(n_preds):
        if _dominant_gt_id_for_pred(pi, matches) != gid:
            continue
        coverage = len(pred_to_gt_hits[(pi, gid)]) / n_gt
        if coverage >= MEANINGFUL_COVERAGE_THRESH:  # 0.30
            meaningful_preds += 1
    if meaningful_preds >= 2:
        n_over += 1
\end{lstlisting}

\subsubsection{Under-Merge (UM) Rate}

An under-merge occurs when a single predicted polyline meaningfully
covers $\geq 2$ distinct GT polylines:

\begin{lstlisting}
# isq_core.py -- compute_isq_image()
for pi, poly in enumerate(pred_polylines):
    n_pred = len(poly)
    meaningful_gts = 0
    for gid in gt_ids:
        n_hit = len(pred_to_gt_hits.get((pi, gid), set()))
        prec_pg = n_hit / n_pred
        rec_pg  = n_hit / max(len(gt_seg_by_id[gid]), 1)
        if prec_pg > 0.3 or rec_pg > 0.3:
            meaningful_gts += 1
    if meaningful_gts >= 2:
        under_merge += 1
\end{lstlisting}

\subsubsection{ISQ Summary Aggregation}

\begin{lstlisting}
# isq_core.py -- summarize_isq()
return {
    "precision":   float(np.mean([m["precision"]    for m in per_image])),
    "recall":      float(np.mean([m["recall"]       for m in per_image])),
    "f1":          float(np.mean([m["f1"]           for m in per_image])),
    "over_split":  float(np.mean([m["over_split_count"] for m in per_image])),
    "under_merge": float(np.mean([m["under_merge"]  for m in per_image])),
    "n_images":    len(per_image),
    # Extended rates (added in eval_isq post-summary):
    # os_rate = OS_total / n_gt_instances
    # um_rate = UM_total / n_pred_polylines_total
}
\end{lstlisting}

\subsubsection{OS and UM Rate Definitions}

In addition to per-image means (\texttt{over\_split}, \texttt{under\_merge}),
we report dataset-level rates used in the presentation (slide~23):

\begin{align*}
  \text{OS rate (\%)} &= 100 \times \frac{N_{\text{OS}}}{\sum_i n_{\text{gt\_instances}}^{(i)}} \\
  \text{UM rate (\%)} &= 100 \times \frac{\sum_i \text{UM}^{(i)}}{\sum_i n_{\text{pred\_polylines}}^{(i)}}
\end{align*}

where $N_{\text{OS}}$ is the total over-split GT-instance count across all images,
and $\text{UM}^{(i)}$ is the under-merge count for image~$i$
(number of predicted polylines that meaningfully cover $\geq 2$ GT instances).
The UM denominator uses \textbf{all} decoded predicted polylines
(including GNN-internal candidates), not only polylines retained for ISQ matching.

\begin{table}[H]
\centering
\caption{ISQ results --- full matrix (slide~23). TTPLA test unless noted.}
\small
\begin{tabular}{llrrrrrr}
\toprule
Resolution & Model & Exp & P & R & F1 & OS (\%) & UM (\%) \\
\midrule
512$^2$  & Geom & exp80 / ep0  & 0.354 & 0.663 & 0.455 & 0.76 & 26.7 \\
512$^2$  & GAT  & exp83 / ep6   & 0.361 & 0.693 & 0.469 & 1.17 & 13.6 \\
1024$^2$ & Geom & exp19          & 0.691 & 0.749 & 0.706 & 0.49 & 18.9 \\
1024$^2$ & GAT  & exp76 / ep5    & 0.770 & 0.788 & 0.769 & 2.65 & 5.2 \\
KEPCO val & Geom & exp23          & 0.436 & 0.320 & 0.336 & 0.40 & 37.6 \\
KEPCO val & GAT  & exp71          & 0.625 & 0.382 & 0.445 & 1.61 & 5.2 \\
\bottomrule
\end{tabular}
\end{table}

\noindent
\textbf{Reproduction.}
512$^2$ rows: \texttt{eval\_isq/results/isq\_results\_exp80\_exp83\_ep6\_512\_test.json}
(220 TTPLA test images, \texttt{YOLinO\_benchmark}).
1024$^2$ rows: \texttt{eval\_isq/results/isq\_results\_exp76\_ep5\_test.json}
(440 TTPLA test images, \texttt{ttpla\_yolino\_dataset\_1024x1024}).
KEPCO rows: \texttt{eval\_isq/results/isq\_results\_kepco\_fitlines.json}
(36 validation frames; geom checkpoint \texttt{exp23\_finetune\_kepco}, GNN \texttt{exp71\_gnn\_kepco}).

At $1024^2$, GAT improves ISQ F1 from 0.706 to 0.769 while reducing UM rate
from 18.9\% to 5.2\% (slide~24 conclusion).
At $512^2$, exp83 improves recall (+3.0\,pp) and F1 (+1.4\,pp) over exp80
at the cost of a slightly higher OS rate (1.17\% vs.\ 0.76\%).

\subsubsection{Motivating Examples (slide~22)}

\paragraph{Line-width invariance.}
Pixel F1 (LSNetv2 GT-B) is sensitive to rasterisation width, whereas ISQ
operates on matched grid segments and is width-invariant.
Representative test tile \texttt{58\_00191\_R} (exp76 geom, $1024^2$):

\begin{table}[H]
\centering
\caption{Line-width sensitivity: pixel F1 vs.\ ISQ (stem \texttt{58\_00191\_R})}
\begin{tabular}{lcc}
\toprule
Metric & @ 2\,px line width & @ 4\,px line width \\
\midrule
Pixel F1 (GT-B) & 0.315 & 0.524 \\
ISQ F1          & \multicolumn{2}{c}{0.897 (identical)} \\
\bottomrule
\end{tabular}
\end{table}

\noindent
Reproduce: \texttt{python eval\_pixel\_f1/viz\_isq\_motivation.py}
(selects stems by scanning \texttt{isq\_results\_exp76\_ep5\_test.json}).

\paragraph{Under-merge case study.}
Tile \texttt{104\_495\_L} shows a merged geom prediction (UM\,=\,3) that
achieves high pixel F1 but low ISQ F1; GAT separates instances (UM\,=\,0):

\begin{table}[H]
\centering
\caption{Under-merge motivating example (stem \texttt{104\_495\_L}, 2\,px GT-B)}
\begin{tabular}{lrrr}
\toprule
Post-processing & Pixel F1 & ISQ F1 & UM count \\
\midrule
Geom (exp19)    & 0.741 & 0.527 & 3 \\
GAT (exp76 ep5) & 0.648 & 0.902 & 0 \\
\bottomrule
\end{tabular}
\end{table}

\subsubsection{ISQ Validation Tests}

Three validation modes are implemented in \texttt{eval\_isq/validate\_isq.py}:

\begin{enumerate}
  \item \textbf{Sanity test}: GT polylines used directly as predictions;
        expected result: precision $= 1.0$, recall $= 1.0$, F1 $= 1.0$,
        OS $= 0$, UM $= 0$.

  \item \textbf{Controlled test}: Artificially generated under-merge
        (concatenate two adjacent GT wires into one polyline) and
        over-split (split a GT wire in half) predictions;
        expected result: UM $> 0$ and OS $> 0$ respectively.

  \item \textbf{Monotonicity test}: As predictions are progressively degraded
        (random jitter, random drop), F1 should decrease monotonically.
\end{enumerate}

%% ============================================================
\section{Inference, Postprocessing, and FPS Measurement}
\label{sec:inference}

\subsection{Stage 1 Postprocessing: BFS + Fit-Lines}
\label{subsec:postproc}

After the geometry head produces per-cell segments, they are assembled into polylines
by the YOLinO rule-based pipeline (\texttt{fit\_lines}):

\begin{enumerate}[leftmargin=*]
  \item \textbf{Adjacency graph}: Segments within squared endpoint distance
        $\leq 768 \,\text{px}^2$ (default) are connected by a directed edge.
  \item \textbf{Connected components}: BFS over the adjacency graph groups
        segments into wire candidates.
  \item \textbf{Spline fitting} (optional): Cubic spline smoothing is applied
        to each connected component.
\end{enumerate}

\begin{lstlisting}
# isq_core.py -- _geom_segments_to_fitlines_cc()
adjacency, reversed_adjacency = get_adjacency_list(
    float(adjacency_threshold), endpoints, startpoints)
roots = [node for node in adjacency.keys()
         if not adjacency[node]]
_, seg_id_groups = get_connected_components(
    confidences, "isq_geom",
    np.zeros((1,1,3), dtype=np.uint8),
    int(min_segments_for_polyline),
    _Args(), reversed_adjacency, roots,
    list(segments), write_debug_images=False)
\end{lstlisting}

\subsection{Stage 2 Postprocessing: GNN Connected Components}
\label{subsec:gnn_postproc}

After GAT edge scoring, predicted polylines are recovered by
\emph{union-find} over edges with probability $\geq$ \texttt{gnn\_cc\_edge\_thresh}:

\begin{lstlisting}
# isq_core.py -- gnn_segments_to_polylines()
node_mid    = _to_numpy(gnn_single["node_mid_px"])
node_valid  = _to_numpy(gnn_single["node_valid"])
edge_scores = _to_numpy(gnn_single["edge_scores"])
neighbors   = _to_numpy(gnn_single["neighbors"])
# Union-Find grouping on thresholded edges
# -> list of segment groups (polylines)
\end{lstlisting}

\subsection{FPS Measurement}
\label{subsec:fps}

FPS is measured in \texttt{src/yolino/tools/benchmark\_yolino\_fps.py}
with three scopes:

\begin{table}[H]
\centering
\caption{FPS benchmark scopes}
\begin{tabular}{lp{9cm}}
\toprule
Scope & Description \\
\midrule
\texttt{geom\_only} & Backbone + FPN + geometry/embed heads only (no postprocessing). \\
\texttt{geom\_post} & Above + activations + grid decode + \texttt{fit\_lines} (with spline). \\
\texttt{geom\_post\_cc} & Forward + conf-filtered tensor decode + adjacency CC only (pre-spline). \\
\texttt{geom\_gnn}  & Full forward including GNN head (\texttt{e2e\_mode=gnn}). \\
\bottomrule
\end{tabular}
\end{table}

\noindent
Presentation slide~19 maps YOLinO\,/\,YOLinO+Post\,/\,YOLinO+GAT to
\texttt{geom\_only} / \texttt{geom\_post\_cc} / \texttt{geom\_gnn}
respectively.
The full \texttt{geom\_post} path (BFS + spline) is $\sim$1.2\,FPS at both
resolutions and is \textbf{not} used in the slide FPS table.

Timing methodology:
\begin{enumerate}[leftmargin=*]
  \item Images are preloaded to GPU memory before the timed loop.
  \item 20 warmup iterations are discarded.
  \item Wall-clock time via \texttt{time.perf\_counter()} with
        \texttt{torch.cuda.synchronize()} before and after each forward pass.
  \item Mean / median / throughput FPS are reported over $N$ images.
\end{enumerate}

\begin{lstlisting}
# benchmark_yolino_fps.py -- _bench_forward()
for i in range(warmup):          # warmup
    fn(batches[i % len(batches)])
_sync(device)                    # torch.cuda.synchronize()
for images in batches:
    _sync(device)
    t0 = time.perf_counter()
    fn(images)
    _sync(device)
    latencies.append(time.perf_counter() - t0)
fps_mean = 1.0 / statistics.mean(latencies)
\end{lstlisting}

FPS measurements are performed at both $512 \times 512$ (benchmark dataset)
and $1024 \times 1024$ (full TTPLA tiles) resolutions on a single GPU
with batch size 1.

\begin{table}[H]
\centering
\caption{Inference FPS comparison (slide~19). 100 val images, batch\,=\,1, CUDA-sync timing.}
\begin{tabular}{lccc}
\toprule
Image size & YOLinO & YOLinO + Post & YOLinO + GAT \\
\midrule
512$^2$  & 131.3 & 54.1 & 44.7 \\
1024$^2$ & 103.4 & 48.0 & 42.3 \\
\bottomrule
\end{tabular}
\end{table}

\noindent
\textbf{Checkpoints and sources.}
512$^2$: \texttt{Final\_comparison/ttpla\_fps\_512\_exp80\_exp83/fps\_summary\_audit.json}
(exp80 geom / exp83 GNN; scopes \texttt{exp80\_geom\_only},
\texttt{exp80\_geom\_post\_cc}, \texttt{exp83\_geom\_gnn}).
1024$^2$ geom-only and GAT: \texttt{Final\_comparison/ttpla\_fps/fps\_summary.json}
(exp19 ep42 geom; exp76 GNN best).
1024$^2$ +Post: \texttt{geom\_post\_cc} scope on exp19
(\texttt{benchmark\_yolino\_fps.py --scope geom\_post\_cc}); reproduced
$\approx$51\,FPS on the evaluation GPU (2026-06-18); slide reports 48.0\,FPS.

\begin{lstlisting}[language=bash, numbers=none]
# 512^2 FPS (matches slide 19)
cd CAPSTONE && CUDA_VISIBLE_DEVICES=0 PYTHONPATH=src YOLINO_IGNORE_DIRTY=1 \
  python src/yolino/tools/benchmark_yolino_fps.py \
    --scope geom_only --config configs/experiments/exp80_ttpla_512512_scale16.yaml \
    --checkpoint ttpla_train_exp/log/checkpoints/exp80_ttpla_512512_scale16/best_model.pth \
    --dataset-root ../YOLinO_benchmark --split val --num-images 100 --gpu

# Repeat with --scope geom_post_cc and exp83 config for +Post / +GAT rows.

# 1024^2 FPS
python src/yolino/tools/benchmark_yolino_fps.py \
  --scope geom_post_cc \
  --config configs/experiments/exp19_fpn_bottomup_p4_num_predictors4_1024.yaml \
  --checkpoint ttpla_train_exp/log/checkpoints/exp19_fpn_bottomup_p4_num_predictors4_1024/ep0042_model.pth \
  --dataset-root ../ttpla_yolino_dataset_1024x1024 --split val --num-images 100 --gpu
\end{lstlisting}

%% ============================================================
\section{Experimental Environment}
\label{sec:environment}

\subsection{Software Stack}

\begin{table}[H]
\centering
\caption{Software environment}
\begin{tabular}{ll}
\toprule
Component & Version / Description \\
\midrule
Python & 3.10+ (project venv at \texttt{/home/work/caps\_drone/yolino/venv}) \\
PyTorch & 2.x (CUDA 11.8+) \\
torchvision & Compatible with PyTorch version above \\
timm & Optional (for HRNet/ResNet backbone variants) \\
NumPy & $\geq$ 1.23 \\
OpenCV (\texttt{cv2}) & $\geq$ 4.6 \\
Shapely & $\geq$ 2.0 (label geometry intersection) \\
TensorBoard & Logging backend \\
huggingface\_hub & Dataset / checkpoint download \\
\bottomrule
\end{tabular}
\end{table}

\subsection{Hardware}

\begin{table}[H]
\centering
\caption{Training and evaluation hardware}
\begin{tabular}{ll}
\toprule
Component & Specification \\
\midrule
GPU (training) & 4 $\times$ NVIDIA GPU (DDP, \texttt{CUDA\_VISIBLE\_DEVICES=0,1,2,3}) \\
GPU (evaluation) & Single GPU (\texttt{CUDA\_VISIBLE\_DEVICES=0}) \\
Distributed training & PyTorch \texttt{torch.distributed.run} (DDP) \\
Mixed precision & Disabled (\texttt{amp: false}) \\
DataLoader workers & 0 (single-process, \texttt{loading\_workers: 0}) \\
\bottomrule
\end{tabular}
\end{table}

\subsection{Reproducibility}

All experiments are launched via \texttt{run.sh}, which sets
\texttt{DATASET\_TTPLA}, \texttt{PYTHONPATH}, and calls
\texttt{torch.distributed.run}:

\begin{lstlisting}[language=bash, numbers=none]
# Stage 1 training
export DATASET_TTPLA="$(pwd)/YOLinO_benchmark"
bash run.sh \
  --config configs/experiments/exp80_ttpla_512512_scale16.yaml \
  --dataset-root "$DATASET_TTPLA" \
  --nproc 4

# Stage 2 training (requires Stage 1 checkpoint)
bash run.sh \
  --config configs/experiments/exp83_gnn_ttpla_512512_from_exp80.yaml \
  --dataset-root "$DATASET_TTPLA" \
  --nproc 3
\end{lstlisting}

Checkpoints are written to:
\begin{lstlisting}[language=bash, numbers=none]
ttpla_train_exp/log/checkpoints/<run_name>/best_model.pth
ttpla_train_exp/log/checkpoints/<run_name>/ep<NNNN>_model.pth
\end{lstlisting}

All report checkpoints are published on Hugging Face
(\texttt{V8heart/CAPSTONE-gnn-weights}).
The \textbf{final Stage~2 model} is \texttt{exp83}, epoch~6
(\texttt{ep0006\_model.pth}); \texttt{exp81/ep0058} is retained as a reference only.

\begin{table}[H]
\centering
\caption{Published checkpoints on Hugging Face (\texttt{CAPSTONE-gnn-weights})}
\begin{tabular}{lll}
\toprule
Experiment & Checkpoint file & Role \\
\midrule
\texttt{exp80\_ttpla\_512512\_scale16} & \texttt{best\_model.pth} & Stage~1 geom (512$^2$) \\
\texttt{exp83\_gnn\_ttpla\_512512\_from\_exp80} & \texttt{ep0006\_model.pth} & \textbf{Stage~2 GNN (final)} \\
\texttt{exp81\_gnn\_ttpla\_512512\_from\_exp80} & \texttt{ep0058\_model.pth} & Previous Stage~2 \\
\texttt{exp19\_fpn\_bottomup\_p4\_num\_predictors4\_1024} & \texttt{ep0042\_model.pth} & 1024$^2$ ablation geom \\
\texttt{exp76\_gnn\_ttpla\_full} & \texttt{ep0005\_model.pth} & 1024$^2$ ablation GNN \\
\texttt{exp23\_finetune\_kepco} & \texttt{best\_model.pth} & KEPCO geom fine-tune \\
\texttt{exp71\_gnn\_kepco} & \texttt{best\_model.pth} & KEPCO GNN \\
\bottomrule
\end{tabular}
\end{table}

\begin{lstlisting}[language=bash, numbers=none]
# Stage 2 final model (exp83, epoch 6)
hf download V8heart/CAPSTONE-gnn-weights \
  exp83_gnn_ttpla_512512_from_exp80/ep0006_model.pth \
  --repo-type model \
  --local-dir ttpla_train_exp/log/checkpoints/exp83_gnn_ttpla_512512_from_exp80

# 1024^2 ablation geom + GNN
hf download V8heart/CAPSTONE-gnn-weights \
  exp19_fpn_bottomup_p4_num_predictors4_1024/ep0042_model.pth \
  --repo-type model \
  --local-dir ttpla_train_exp/log/checkpoints/exp19_fpn_bottomup_p4_num_predictors4_1024
hf download V8heart/CAPSTONE-gnn-weights \
  exp76_gnn_ttpla_full/ep0005_model.pth \
  --repo-type model \
  --local-dir ttpla_train_exp/log/checkpoints/exp76_gnn_ttpla_full

# KEPCO cross-domain
hf download V8heart/CAPSTONE-gnn-weights \
  exp23_finetune_kepco/best_model.pth \
  --repo-type model \
  --local-dir ttpla_train_exp/log/checkpoints/exp23_finetune_kepco
hf download V8heart/CAPSTONE-gnn-weights \
  exp71_gnn_kepco/best_model.pth \
  --repo-type model \
  --local-dir ttpla_train_exp/log/checkpoints/exp71_gnn_kepco
\end{lstlisting}

Published datasets:
\begin{itemize}[leftmargin=*]
  \item \texttt{V8heart/yolino-ttpla-benchmark} --- 512$^2$ main benchmark
  \item \texttt{V8heart/ttpla-yolino-1024} --- 1024$^2$ ablation / ISQ test tiles
  \item \texttt{V8heart/Yolino-KEPCO} --- KEPCO cross-domain evaluation set
\end{itemize}

%% ============================================================
\section*{References}
\addcontentsline{toc}{section}{References}

\begin{thebibliography}{9}

\bibitem{meyer2021yolino}
A.\ Meyer, P.\ Skudlik, J.-H.\ Pauls, and C.\ Stiller,
``YOLinO: Generic Single Shot Polyline Detection in Real Time,''
\textit{Proc.\ IEEE/CVF ICCV Workshops}, pp.\ 2916--2925, 2021.

\bibitem{velickovic2018gat}
P.\ Veli\v{c}kovi\'{c}, G.\ Cucurull, A.\ Casanova, A.\ Romero,
P.\ Li\`{o}, and Y.\ Bengio,
``Graph Attention Networks,''
\textit{ICLR}, 2018.

\bibitem{liu2022convnext}
Z.\ Liu, H.\ Mao, C.-Y.\ Wu, C.\ Feichtenhofer, T.\ Darrell, and S.\ Xie,
``A ConvNet for the 2020s,''
\textit{Proc.\ IEEE/CVF CVPR}, pp.\ 11976--11986, 2022.

\bibitem{lin2017fpn}
T.-Y.\ Lin, P.\ Doll\'{a}r, R.\ Girshick, K.\ He, B.\ Hariharan, and S.\ Belongie,
``Feature Pyramid Networks for Object Detection,''
\textit{Proc.\ IEEE/CVF CVPR}, pp.\ 2117--2125, 2017.

\bibitem{liu2018panet}
S.\ Liu, L.\ Qi, H.\ Qin, J.\ Shi, and J.\ Jia,
``Path Aggregation Network for Instance Segmentation,''
\textit{Proc.\ IEEE/CVF CVPR}, pp.\ 8759--8768, 2018.

\bibitem{he2015msra}
K.\ He, X.\ Zhang, S.\ Ren, and J.\ Sun,
``Delving Deep into Rectifiers: Surpassing Human-Level Performance on ImageNet Classification,''
\textit{Proc.\ IEEE ICCV}, pp.\ 1026--1034, 2015.

\bibitem{ttpla}
Q.\ Wang and J.\ Zhang,
``TTPLA: An Aerial-Image Dataset for Detection and Segmentation of Transmission Towers and Power Lines,''
\textit{Proc.\ ACM MM}, 2020.

\bibitem{lsnetv2}
X.\ Li \textit{et al.},
``LSNetv2: Powerline Detection via Feature Line Segments,''
2023.

\end{thebibliography}

\end{document}
